% Gerado por tcc/analise/resultados_c3.py -- não edite à mão; rode o script de novo.
\begin{table}[htb]
\caption{Diferença entre os grupos, intervalo de confiança e tamanho de efeito}\label{tab:efeito}
\centering\footnotesize\setlength{\tabcolsep}{4pt}
\begin{tabular}{llrcc}
\toprule
Indicador & Período & F $-$ NF & IC 95\% & g de Hedges \\
\midrule
Dívida bruta / Ativo total & 2018--2025 & -4,5 p.p. & [-21,3; 12,2] & -0,32 \\
Exigível total / Ativo total & 2018--2025 & -0,8 p.p. & [-13,0; 11,4] & -0,07 \\
Dívida bruta / (Dívida bruta + PL) & 2018--2025 & -6,3 p.p. & [-30,6; 18,0] & -0,33 \\
Dívida líquida / EBITDA & 2018--2025 & -3,0x & [-9,3; 3,4] & -0,76 \\
Dívida CP / Dívida total & 2018--2025 & -17,8 p.p. & [-41,7; 6,1] & -1,28 \\
Dívida CP / Dívida total & 2020--2021 & -16,8 p.p. & [-32,6; -1,0] & -1,39 \\
\bottomrule
\end{tabular}
\fonte{Elaborado pelo autor (2026), com dados da Economática.}
\nota{Diferença entre as médias das familiares e das não familiares (p.p.: pontos percentuais), com intervalo de confiança de 95\% pela aproximação de Welch. g de Hedges: diferença padronizada, com correção para amostras pequenas; por convenção, valores em torno de 0,2, 0,5 e 0,8 indicam efeitos pequeno, médio e grande.}
\end{table}
